\SourceLeaf{381}%
% A544-P0381-D01
\[
\left.\begin{aligned}
f&\equiv i^{-\mathrm{P}\varphi}\varphi\\
i^{-\mathrm{P}mf}mf&\equiv i^{-\mathrm{P}m\varphi}m\varphi
\end{aligned}\right\}\ \text{mod. }\mu
\]
woraus leicht folgt $\mathrm{P}mf-\mathrm{P}m\varphi\equiv-\mathrm{P}\varphi\ (\text{mod. }4)$ und das Aggregat der beiden obersten Reihen $\equiv-\Sigma\mathrm{P}\varphi$. Da nun ferner $\mathrm{P}m\varphi-\Pi m\varphi\equiv\mathrm{P}(i^{-\Pi m\varphi}m\varphi)$ ist, $m\varphi\pdot i^{-\Pi m\varphi}$ zu $\kappa$ gehört und der Inbegriff \emph{aller} $m\varphi\pdot i^{-\Pi m\varphi}$ ohne Rücksicht auf die Ordnung mit allen $\varphi$ übereinkommt, so wird das Aggregat \emph{aller} $\mathrm{P}(m\varphi\pdot i^{-\Pi m\varphi})$ aequal sein dem Aggregat aller $\mathrm{P}\varphi$; folglich das Aggregat der dritten und vierten Reihe $\equiv\Sigma\mathrm{P}\varphi\ (\text{mod. }4)$, also das Aggregat aller vier Reihen $\equiv0\ (\text{mod. }4)$ W.\ Z.\ B.\ W.

% A544-P0381-D02
Da also $\varepsilon$, unabhängig von der Wahl der Viertelsysteme blos von $m$ und $\mu$ abhängt, so werden wir $\varepsilon$ den Character der Zahl $m$ in Beziehung auf den Modulus $\mu$ nennen und mit $\text{Ch. }m\ (\text{mod. }\mu)$ bezeichnen. Man sieht leicht, dass dies nur eine Generalisirung derjenigen Definition ist, die (Art.\ $\ldots$) für den Fall, wo $\mu$ eine Primzahl ist, gegeben ist, oder sie unter sich begreift.

% A544-P0381-D03
\Article{4}
Ich gehe jetzt zu \emph{bestimmten} Anordnungen der Viertelsysteme über, und werde den mit $m$ zu bezeichnenden Modulus $=ea+ebi$ setzen, so dass die positive ganze Zahl $e$ den grössten reellen Divisor, oder den grössten Divisor, welchen die beiden Bestandtheile von $m$ haben, bedeutet, oder $a,b$ Primzahlen unter sich. Das am einfachsten angeordnete Viertelsystem wird das sein, dessen Glieder $x+iy$ so beschaffen sind, dass $ax+by$ positiv und kleiner als $\frac12e(aa+bb)$, $ay-bx$ nicht negativ, und gleichfalls kleiner als $\frac12e(aa+bb)$ wird; die letztere Bedingung wird geflissentlich so ausgedrückt, dass auch die Fälle, wo $ay-bx=0$ wird, darunter begriffen sind. Man sieht leicht, dass solcher Fälle zusammen $\frac12(e-1)$ sein werden, nemlich

% A544-P0381-D04
\[
\begin{array}{rlrl}
x&=a,&y&=b\\
x&=2a,&y&=2b\\
x&=3a,&y&=3b\\
\multicolumn{4}{c}{\text{u.\,s.\,w. bis}}\\
x&=\frac12(e-1)a,&y&=\frac12(e-1)b
\end{array}
\]
also gar keine, wenn die Bestandtheile von $m$ keinen gemeinschaftlichen Divisor %
\SourceLeaf{382}%
% A544-P0382-D01
haben. Nennen wir dieses Viertelsystem $k$, und $k',k'',k'''$ diejenigen, welche entstehen, indem man die zu $k$ gehörigen Zahlen mit $i,-1,-i$ multiplicirt, oder man mag auch setzen
\[
k'=m+ik,\quad k''=(1+i)m-k,\quad k'''=im-ik
\]

% A544-P0382-D02
Auf diese Art erhält man folgende Regel, um zu beurtheilen, ob eine beliebige vorgegebene durch $m$ nicht theilbare ganze Zahl $x+iy$ congruent sei einem Gliede von $k,k',k''$ oder $k'''$, nemlich indem man kann $2(ax+by)$ in die Form $Pe(aa+bb)+Q$, $2(ay-bx)$ in die Form $Re(aa+bb)+S$ bringen, so dass $P,Q,R,S$ ganze reelle Zahlen und zwar

% A544-P0382-D03
\begin{center}wenn $x+iy$ congruent ist einer Zahl aus\end{center}
\begingroup\small\setlength{\arraycolsep}{3pt}
\[
\begin{array}{l||l|l|l|l}
\text{so dass}&k&k'&k''&k'''\\
P&\text{gerade}&\text{ungerade}&\text{ungerade}&\text{gerade}\\
R&\text{gerade}&\text{gerade}&\text{ungerade}&\text{ungerade}\\
Q&\text{positiv}&\text{positiv}&\text{positiv}&\text{nicht negativ}\\
e(aa+bb)-Q&\text{positiv}&\text{nicht negativ}&\text{positiv}&\text{positiv}\\
S&\text{nicht negativ}&\text{positiv}&\text{positiv}&\text{positiv}\\
e(aa+bb)-S&\text{positiv}&\text{positiv}&\text{nicht negativ}&\text{positiv}
\end{array}
\]
\endgroup

% A544-P0382-D04
Man erleichtert sich die Uebersicht, wenn man die Fälle, wo keine der Zahlen $ax+by$, $ay-bx$ durch $e(aa+bb)$ theilbar ist, von den übrigen unterscheidet.

I. Im ersten Falle hat man für $P$ schlechthin die (algebraisch) kleinere der beiden ganzen Zahlen zu nehmen, zwischen welche (ausschliesslich) $\frac{2(ax+by)}{e(aa+bb)}$ fallen wird, und eben so für $R$ die kleinere der beiden, zwischen welche $\frac{2(ay-bx)}{e(aa+bb)}$ fällt.

% A544-P0382-D05
II. Ist $ax+by$ durch $e(aa+bb)$ theilbar, so wird $ay-bx$ zwar durch $aa+bb$, nicht aber durch $e(aa+bb)$ theilbar sein (weil sonst $x+iy$ durch $ea+ebi$ theilbar sein würde). Ist nun $R$, d.\,i. die Zahl, welche zunächst kleiner ist als $\frac{2(ay-bx)}{e(aa+bb)}$, gerade, so wird $x+iy$ einer Zahl aus $k'$ congruent sein nach dem Mod.\ $ea+ebi$, einer aus $k'''$ hingegen, wenn $R$ ungerade ist.

% A544-P0382-D06
III. Ist $ay+bx$ durch $e(aa+bb)$ theilbar, nicht aber $ax+by$, so wird $x+iy$ einer Zahl aus $k$, oder aus $k''$ congruent sein, je nachdem die ganze Zahl, welche algebraisch zunächst kleiner ist als $\frac{2(ax+by)}{e(aa+bb)}$, gerade oder ungerade ist.

\SourceLeaf{383}%
% A544-P0383-D01
\Article{5}
Man leitet aus obigem ohne Mühe folgende Methode ab zur Bestimmung des Characters einer gegebenen ganzen Zahl $M=A+Bi$ in Beziehung auf den ungeraden sie nicht messenden Modulus $m=ea+ebi$.

Zur Abkürzung bedienen wir uns folgender Bezeichnung. Wenn $p$ irgend eine gebrochene reelle Zahl vorstellt, soll durch $[p]$ diejenige ganze Zahl bezeichnet werden, die zugleich $p-[p]$ und $1+[p]-p$ positiv macht. Bei dieser Definition ist also die Anwendung der Bezeichnung auf ganze Zahlen ausgeschlossen. Fasste man die Definition so, dass weder $p-[p]$ noch $1+[p]-p$ negativ sein soll, so würde das Zeichen $[p]$ für den Fall, wo $p$ ganze Zahl ist, zweideutig sein. Man könnte auch, wie in einer früheren Abhandlung geschehen ist, die Bedingung so stellen, dass $1+[p]-p$ positiv und $p-[p]$ nur nicht negativ sein soll. Für unsern Zweck ist es etwas bequemer, sich an die erste Begriffsbestimmung zu halten.

% A544-P0383-D02
Das Viertelsystem $k$ bilden hienach alle ganzen Zahlen $f$, wofür wenn man $\frac{2f}{m}=\xi+\eta i$ setzt, $\xi$ zwischen $0$ und $1$ ausschliesslich, $\eta$ zwischen $0$ und $1$, die $0$ eingeschlossen liegt, oder
\[
\begin{array}{lll}
&[\xi]=0,&[\eta]=0\\
\text{oder}&[\xi]=0,&\eta=0
\end{array}
\]

% A544-P0383-D03
Setzt man nun für jedes $f$, $\frac{2fM}{m}=\xi+i\eta$ und nimmt
\begingroup\small\setlength{\arraycolsep}{2pt}
\[
\begin{array}{rl}
\Psi f=0&\text{wenn zugleich }[\xi]\text{ gerade und entweder }[\eta]\text{ gerade oder }\eta\text{ ganz}\\
1&[\eta]\text{ gerade und entweder }[\xi]\text{ ungerade oder }\xi\text{ ganz}\\
2&[\xi]\text{ ungerade und entweder }[\eta]\text{ ungerade oder }\eta\text{ ganz}\\
3&[\eta]\text{ ungerade und entweder }[\xi]\text{ gerade oder }\xi\text{ ganz}
\end{array}
\]
\endgroup

% A544-P0383-D04
tabellarisch so
\[
\begin{array}{l|c|c|c}
&[\eta]\text{ gerade}&[\eta]\text{ ungerade}&\eta\text{ ganz}\\\hline
[\xi]\text{ gerade}&0&3&0\\
\relax[\xi]\text{ ungerade}&1&2&2\\
\xi\text{ ganz}&1&3&-\ \pdot
\end{array}
\]
was man durch $\nabla(\xi+i\eta)$ bezeichnen mag, so wird der gesuchte Character der Zahl $M$ in Beziehung auf den Modulus $m$ aequal dem kleinsten Reste von $\Sigma\Psi f$ nach dem Modulus $4$.

\SourceLeaf{384}%
% A544-P0384-D01
\Article{6}
Die im vorhergehender Art.\ gegebene Vorschrift ist allgemein: für den Fall, wo $M$ ungerade ist, werden wir ihr aber eine andere Gestalt geben. Wir werden zugleich annehmen, dass die reellen Theile von $m$ und $M$ ungerade, also die imaginären gerade sind.

% A544-P0384-D02
Wir lassen jeder zu dem Viertelsysteme $k$ gehörenden Zahl $f$ eine andere $g$ correspondiren, die aus $f$ auf folgende Art abgeleitet wird. Indem man $\frac{2f}{m}=\xi+i\eta$ setzt, unterscheidet man vier Fälle
\[
\begin{array}{ll}
\mathrm{I.}&\text{Wenn }[2\xi]=0\text{ und entweder }[2\eta]=0\text{ oder }\eta=0\\
\mathrm{II.}&\text{Wenn }[2\xi]=1\text{ und entweder }[2\eta]=0\text{ oder }\eta=0\\
\mathrm{III.}&\text{Wenn }[2\xi]=0\text{ und }[2\eta]=1\\
\mathrm{IV.}&\text{Wenn }[2\xi]=1\text{ und }[2\eta]=1
\end{array}
\]
Im ersten Falle wird man $g=2f$, im zweiten $g=im-2if$, im dritten $g=m+2if$, im vierten $g=(1+i)m-2f$ setzen. Man sieht leicht, dass der Inbegriff aller $g$ ein vollständiges Viertelsystem $l$ bildet; ihre Charakteristik ist, dass zugleich, wenn man $\frac gm=\xi+i\eta$ setzt

% A544-P0384-D03
\[
\begin{array}{llll}
\text{entweder}&[\xi]=0,&[\eta]=0&\\
\text{oder}&\eta=0,&[\xi]=0&\text{und }g\text{ durch }1+i\text{ theilbar}\\
\text{oder}&\xi=0,&[\eta]=0&\text{und }g\text{ durch }1+i\text{ nicht theilbar}
\end{array}
\]
Daraus folgt, dass $l$ sich von $k$ nur dadurch unterscheidet, dass diejenigen Zahlen in $k$, für welche $\eta=0$, und die durch $1+i$ untheilbar sind, nemlich
\[
\begin{aligned}
&a+bi,\quad3(a+bi),\quad5(a+bi)\ldots\frac{e-3}{2}(a+bi)\\
&\text{oder bis }\frac{e-1}{2}(a+bi)
\end{aligned}
\]
je nachdem $e$ von der Form $4n+1$ oder $4n+3$, in $l$ fehlen und dagegen in letzterm Complex die Producte jener Zahlen in $i$ auftreten. Zugleich sieht man, dass für $e=1$, d.\,i. wenn $m$ durch keine reelle ganze Zahl theilbar ist, $k$ und $l$ ganz gleich sind.

% A544-P0384-D04
Es kommt nun darauf an, $\Psi f$ unmittelbar aus dem dem $f$ entsprechenden $g$ abzuleiten. Das Resultat ist, dass für obige $4$ Fälle
\[
\begin{aligned}
\mathrm{I.}\quad&\Psi f=\nabla\frac{gM}{m}\\
\mathrm{II.}\quad&\Psi f=\left\{\begin{array}{ll}
-\nabla\frac{gM}{m}&\begin{gathered}\text{wenn weder reeller noch imaginärer Th.}\\\text{ von }\frac{gM}{m}\text{ ganz}\end{gathered}\\
1-\nabla\frac{gM}{m}&\text{wenn einer von beiden ganz}
\end{array}\right.
\end{aligned}
\]
III. wie II.\quad IV. wie I.
\BookRule{.3\textwidth}

\SourceLeaf{385}%
% A544-P0385-D01
\par\bigskip\Needspace{8\baselineskip}
\pdfbookmark[1]{Bemerkungen zur Theorie der biquadratischen Reste}{nachlass-biquadratic-bemerkungen}
\begin{center}\large\scshape Bemerkungen.\end{center}
\BookRule{.2\textwidth}
Die Bruchstücke, die hier im Druck mit I und II bezeichnet sind, gehören nach dem Orte zu urtheilen, den die betreffenden Handschriften in einem Notizbuche einnehmen, dem Jahre $1811$ oder der zunächst folgenden Zeit an. Von den vorangehenden Versuchen, den Beweis des Fundamentaltheorems für biquadratische Reste nach den hier für den Rest $1+i$ angewandten Methoden durchzuführen, ist eine Aufzeichnung vorhanden, welche den speciellen Fall des Restes $1+2i$ erledigt und von derjenigen Bestimmung des biquadratischen Characters ausgeht, die man als Note dem Art.\ $6$ des Bruchstücks III beigefügt hat. Im übrigen lassen sich die historischen Angaben, die \textsc{Gauss} in den Anzeigen seiner arithmetischen Abhandlungen veröffentlicht hat, mit Hülfe des Nachlasses dahin ergänzen, dass die in den Artt.\ $15$ bis $20$ der \emph{Theoria residuorum biquadrat.} aufgenommenen Lehrsätze schon vor der Ausarbeitung der \emph{Theoria motus corporum coel.} niedergeschrieben sind. Die in den Anzeigen erwähnten Untersuchungen über cubische Reste werden wohl nicht zur Ausarbeitung gelangt sein; aufgezeichnet finden sich davon die mit den Hülfsmitteln, welche die Abhandlung \emph{Disquisitionum circa aequationes puras ulterior evolutio} bietet, durchgeführten Beweise der Reciprocitätssätze für zwei Primzahlen, von denen die eine reell ist.

% A544-P0385-D02
Die Bruchstücke III bis VI bilden in der Handschrift besondere Hefte und für die drei ersten derselben weist die Form der Schriftzüge auf eine Zeit, die der für die Bruchstücke I und II nicht fern liegt, während für das letzte, Nr.\ VI, ein bedeutend späterer Zeitpunkt angenommen werden muss.

% A544-P0385-D03
\par\medskip
[I.] Art.\ $10$. Die Bestimmung der Anzahl der Ganzenpunkte in $(z,z',z'+\zeta,z+\zeta)$ ergiebt sich aus dem Satze: bedeuten $a$ und $b$ relative Primzahlen, so geht die von $\xi+\eta i$ nach $\xi+\eta i+a+bi$ gezogene Gerade durch Einen Ganzenpunkt, wenn der imaginäre Theil von $(\xi+\eta i)\pdot(-a+bi)$ eine ganze Zahl ist.

